\begin{table}[htbp]
\centering
\caption{Result 2: Evidence of License-Cap Policy, Capped vs.\ Uncapped States}
\label{tab:result2}
\small
\begin{tabular}{lcc}
\toprule
Outcome & Coefficient (border effect) & $p$-value (county-cl.) \\
\midrule
\multicolumn{3}{l}{\textit{Colorado + Washington (license caps bind)}} \\
log(revenue per dispensary) & $+0.341$ & $0.033$ \\
log(revenue per capita) & $+0.046$ & $0.833$ \\
\midrule
\multicolumn{3}{l}{\textit{Montana (no statewide numerical cap)}} \\
log(revenue per capita), 2022--2025 panel & $+0.654$ & $0.036$ \\
log(revenue per dispensary), 2025 cross-section & $+0.012$ & $0.958^{a}$ \\
\bottomrule
\end{tabular}
\begin{minipage}{0.95\textwidth}
\vspace{0.3em}
\footnotesize
\textit{Notes:} $^{a}$HC1 standard error; too few Montana border counties ($n=23$) for county-clustered
inference. Colorado and Washington constrain entry (Washington through a statewide license cap allocated by
lottery; Colorado through local governments' authority under Amendment 64 to ban or limit retail marijuana).
Montana imposes no equivalent statewide cap. Where entry is capped, the border effect appears in revenue per
dispensary and not per capita; where entry is not capped, it appears in revenue per capita and not per
dispensary, consistent with a license-cap explanation for why the border premium in Result 1 shows up on the
intensive rather than the extensive margin.
\end{minipage}
\end{table}
